\section{Experimental Design}
\label{sec:experimental}

\subsection{Dataset Collection and Curation}
Our dataset is the cornerstone of the CMM. We collected data in three primary categories to ensure robustness:
\begin{itemize}
    \item \textbf{Positives:} These images contain authentic UNCC merchandise, serving as the ground truth for our target class. We sourced these images from official university photostreams, social media channels, and manual data collection at the campus bookstore. The positive samples include a wide variety of merchandise types, such as hoodies, t-shirts, and hats, featuring different logo variations (e.g., the "Pick" logo, the crown logo, and block "UNC Charlotte" text) to ensure the model learns to generalize across different official designs.
    \item \textbf{Negatives:} To prevent false positives, we curated a set of negative images containing people without any merchandise or with generic, unbranded clothing. These samples are critical for teaching the model to distinguish between a person and a person wearing specific university gear. We included images of crowds, students in plain clothes, and people in other settings to reinforce the model's ability to reject non-relevant subjects.
    \item \textbf{Hard Negatives:} A crucial part of our dataset consists of "hard negatives"—images that closely resemble UNCC gear but are not authentic. This includes clothing with similar color palettes (e.g., green shirts with white text) or generic "Charlotte" city merchandise that is not university-affiliated. By training on these samples, the model learns to discriminate based on specific logo features rather than just color or general text shape, significantly reducing false detections in a city where green and white are common colors.
    \item \textbf{Augmented Samples:} To improve model generalization and robustness against real-world variability, we expanded our dataset from 722 to 1227 images using various augmentation techniques. These included random horizontal flips, small-scale translations, color jitter (HSV), mosaic augmentation, mixup (light), and moderate scaling. This diversity ensures the model remains effective even when the input video feed is imperfect or the subjects are viewed from oblique angles.
\end{itemize}

The final dataset was split into 82\% training, 12\% validation, and 6\% testing sets to ensure rigorous evaluation.

\begin{figure}[h!]
    \centering
    \begin{subfigure}{0.45\linewidth}
        \includegraphics[width=\linewidth]{pos1.png}
        \caption{Positive Samples}
        \label{fig:pos1}
    \end{subfigure}
    \hfill
    \begin{subfigure}{0.45\linewidth}
        \includegraphics[width=\linewidth]{neg1.png}
        \caption{Negative Samples}
        \label{fig:neg1}
    \end{subfigure}
    \vskip\baselineskip
    \begin{subfigure}{0.45\linewidth}
        \includegraphics[width=\linewidth]{hardneg1.png}
        \caption{Hard Negative Samples}
        \label{fig:hardneg1}
    \end{subfigure}
    \hfill
    \begin{subfigure}{0.45\linewidth}
        \includegraphics[width=\linewidth]{spec1.png}
        \caption{Augmented Samples}
        \label{fig:spec1}
    \end{subfigure}
    \caption{Dataset samples demonstrating the four categories used for training: (a) Positive samples with authentic merchandise, (b) Negative samples with generic clothing, (c) Hard Negative samples with similar colors/text, and (d) Augmented samples with noise and rotation.}
    \label{fig:dataset_samples_all}
\end{figure}

\subsection{Model Selection}
We chose YOLOv11 as the base architecture for the Charlotte Merchandise Model (CMM) because it balances real-time performance and improved feature extraction for small and occluded objects, which are common in crowd scenes. YOLOv11 maintains the single-shot detection paradigm that enables low-latency inference, a key requirement for our "School Spirit Meter" application \cite{khanam2024}.

\begin{figure*}[t]
    \centering
    \includegraphics[width=0.8\textwidth]{yolov11_frame.png}
    \caption{YOLOv11 feature frame and block diagram highlighting C3k2, SPPF, and C2PSA components (\textit{Image credit: \cite{yang2025}}).}
    \label{fig:yolov11_frame}
\end{figure*}

YOLOv11 introduces several architectural changes that improve its accuracy and efficiency. Notably, it introduces three components that distinguish it from previous YOLO versions: the C3k2 block, SPPF (Spatial Pyramid Pooling - Fast), and C2PSA (Convolutional block with Parallel Spatial Attention). The following paragraph describes these components in more detail and explains why they are beneficial for our task.

The C3k2 block replaces previous C2f/C3 blocks and is more computationally efficient while preserving representational power; this contributes to higher processing speed without large accuracy trade-offs. The SPPF module accelerates multi-scale feature aggregation, which helps the detector capture logo features across scales. The C2PSA block adds parallel spatial attention that improves localization and discrimination of small or partially occluded logos. Together these innovations make YOLOv11 a good fit for real-time detection of small or partially occluded university logos in crowded scenes \cite{yang2025}.

\begin{figure}[b]
    \centering
    \includegraphics[width=0.9\linewidth]{yolov11.png}
    \caption{YOLOv11 architecture overview (\textit{Image credit: \cite{khanam2024}}).}
    \label{fig:yolov11}
\end{figure}

While other architectures were considered, they were deemed less suitable for our specific constraints. Faster R-CNN \cite{ren2015faster}, a two-stage detector, offers high accuracy but suffers from higher latency, making it impractical for our real-time "School Spirit Meter" which requires immediate feedback on live video feeds. Similarly, while SSD (Single Shot MultiBox Detector) \cite{liu2016ssd} is faster than R-CNN, it often struggles with small object detection compared to modern YOLO iterations. We also evaluated the recent RT-DETR \cite{zhao2024detrs}, a transformer-based real-time detector. Although RT-DETR shows promising results, its computational overhead on edge devices can be higher than CNN-based approaches like YOLOv11, and the latter's specific optimizations for small object detection (via C2PSA) made it the superior choice for detecting university logos in crowded environments.

\subsection{Model Architecture and Training}
We utilized the YOLOv11m model, a state-of-the-art iteration of the YOLO family. The model was initialized with pretrained weights to leverage transfer learning. We fine-tuned the model on our custom dataset using the following hyperparameters:
\begin{itemize}
    \item \textbf{Epochs:} 50
    \item \textbf{Image Size:} 512x512
    \item \textbf{Batch Size:} 16
\end{itemize}
Training was conducted using the Ultralytics framework. The model was trained to detect three specific classes: `UNCC-LOGO`, `UNCC TORSO`, and `UNCC HEADGEAR`.

\subsection{Detection Pipeline}
To accurately count "merchandise" rather than just logos, we implemented a post-processing logic. A valid "merch" detection is defined as a `UNCC-LOGO` bounding box that is spatially contained within a `UNCC TORSO` or `UNCC HEADGEAR` bounding box. This spatial relationship check ensures that we are counting people wearing the merchandise, rather than loose items or false positives. The system runs in real-time, drawing bounding boxes and updating a live "UNCC MERCH COUNT" on the video feed.
